\chapter{Magnetismo e campos magnéticos}\label{ch:g11:magnetic-fields}

Os navegadores se guiaram por agulhas imantadas durante mil anos antes
que alguém soubesse por que elas apontam para o norte. A resposta chegou
em 1820, quando uma corrente fez uma bússola estremecer: o magnetismo é
feito por cargas em movimento. Este capítulo mapeia o \omterm{def:g11:magnetic-fields:bfield}{campo magnético} ---
o dos \omterm{def:g11:magnetic-fields:magnet}{ímãs}, o da Terra, o das correntes --- e depois fecha o ciclo: o
\omterm{def:g11:electric-gravitational-fields:field}{campo} empurra de volta as correntes, e é assim que todo motor elétrico
gira.

\section{Ímãs e o campo magnético}

\begin{definition}[Ímã e polos]\label{def:g11:magnetic-fields:magnet}
Um \emph{ímã}\index{ímã} atrai o ferro e, pendurado por um fio, gira até
encarar o norte. A sua ação se concentra em dois \emph{polos
magnéticos}\index{polo magnético}, \emph{norte} (a ponta que procura o
norte geográfico) e \emph{sul}: polos iguais se repelem, polos diferentes
se atraem. Serrar um ímã ao meio produz dois ímãs completos ---
experiência alguma jamais isolou um polo.
\end{definition}

\begin{definition}[Campo magnético]\label{def:g11:magnetic-fields:bfield}
Um \omterm{def:g11:magnetic-fields:magnet}{ímã} modifica o espaço à sua volta
(\cref{ch:g11:electric-gravitational-fields}): ele cria em toda parte um
\emph{campo magnético}\index{campo magnético} $\vect B$. Uma bússola lê a
direção dele (a agulha se alinha com $\vect B$, com a ponta norte à
frente); uma \emph{sonda Hall}\index{sonda Hall} mede o seu módulo, em
\emph{teslas}\index{tesla} (\unit{T}).
\end{definition}

\begin{definition}[Linhas de campo magnético]\label{def:g11:magnetic-fields:lines}
As \emph{linhas de campo magnético}\index{linha de campo magnético} são
curvas tangentes a $\vect B$, que se apertam onde o \omterm{def:g11:electric-gravitational-fields:field}{campo} é forte. Fora
de um \omterm{def:g11:magnetic-fields:magnet}{ímã}, elas correm do polo norte para o sul; ao contrário das linhas
de \omterm{def:g11:electric-gravitational-fields:efield}{campo elétrico}, elas não têm extremidades: cada uma se fecha através
do corpo do \omterm{def:g11:magnetic-fields:magnet}{ímã}.
\end{definition}

\begin{omfigure}
\begin{tikzpicture}[scale=0.9,
  fline/.style={omThm, postaction={decorate},
    decoration={markings, mark=at position 0.5 with {\arrow{Stealth}}}}]
  \fill[omProp!15] (-1.4,-0.35) rectangle (0,0.35);
  \fill[omDef!15] (0,-0.35) rectangle (1.4,0.35);
  \draw[thick] (-1.4,-0.35) rectangle (1.4,0.35);
  \draw (0,-0.35) -- (0,0.35);
  \node[font=\small] at (-0.7,0) {S};
  \node[font=\small] at (0.7,0) {N};
  \draw[fline] (1.4,0.2) .. controls (2.6,0.9) and (-2.6,0.9) .. (-1.4,0.2);
  \draw[fline] (1.4,0.3) .. controls (3.8,2.0) and (-3.8,2.0) .. (-1.4,0.3);
  \draw[fline] (1.4,-0.2) .. controls (2.6,-0.9) and (-2.6,-0.9) .. (-1.4,-0.2);
  \draw[fline] (1.4,-0.3) .. controls (3.8,-2.0) and (-3.8,-2.0) .. (-1.4,-0.3);
  \draw[fline] (1.4,0) -- (2.95,0);
  \draw[omExo, thick] (0,1.1) circle (0.22);
  \draw[-Stealth, omDef, thick] (0.22,1.1) -- (-0.22,1.1);
  \draw[omExo, thick] (2.45,0) circle (0.22);
  \draw[-Stealth, omDef, thick] (2.23,0) -- (2.67,0);
\end{tikzpicture}

{\small \omterm{def:g11:electric-gravitational-fields:lines}{Linhas de campo} de um \omterm{def:g11:magnetic-fields:magnet}{ímã} em barra: saem do N, entram no S e se
fecham através do corpo. Uma bússola colocada em qualquer ponto se
acomoda ao longo da linha que passa por ele.}
\end{omfigure}

\section{A Terra é um ímã}

\begin{definition}[O campo magnético da Terra]\label{def:g11:magnetic-fields:earth}
Correntes no núcleo de ferro fundido da Terra fazem do planeta um \omterm{def:g11:magnetic-fields:magnet}{ímã}
gigante: na superfície $B \approx \qty{50}{\micro T}$, e uma agulha livre
se alinha com a componente horizontal do \omterm{def:g11:electric-gravitational-fields:field}{campo} --- é a bússola. Os polos
geomagnéticos não são os geográficos; o ângulo entre o norte da bússola e
o norte verdadeiro é a \emph{declinação}\index{declinação}, corrigida em
todo rumo.
\end{definition}

\begin{example}[Corrigir um rumo]\label{ex:g11:magnetic-fields:bearing}
Onde a \omterm{def:g11:magnetic-fields:earth}{declinação} é de $\ang{7}$ para oeste, um navio que siga o ``norte
da bússola'' por \qty{5.0}{km} deriva
$5000 \times \sin\ang{7} \approx \qty{6.1e2}{m}$ a oeste da sua rota:
para navegar rumo ao norte verdadeiro, siga $\ang{7}$ a \emph{leste} da
agulha.
\end{example}

\section{O campo magnético de uma corrente}

\begin{remark}[Oersted, 1820]\label{rem:g11:magnetic-fields:oersted}
Durante uma aula em 1820, Hans Christian Oersted viu a agulha de uma
bússola girar quando fechou um circuito ali perto: eletricidade e
magnetismo eram um único assunto --- \emph{as correntes criam \omterm{def:g11:magnetic-fields:bfield}{campos
magnéticos}}.
\end{remark}

\begin{proposition}[Campo de um fio retilíneo]\label{prop:g11:magnetic-fields:wire}
Um fio retilíneo longo percorrido por uma corrente $I$ cria um \omterm{def:g11:electric-gravitational-fields:field}{campo}
cujas linhas são círculos centrados no fio, em planos perpendiculares a
ele. À distância $d$ do fio,
\[
B = \frac{\mu_0\, I}{2\pi d},
\qquad \mu_0 = 4\pi\times 10^{-7}\ \unit{T.m/A} \approx \qty{1.26e-6}{T.m/A}.
\]
\end{proposition}
\admitted

\begin{omfigure}
\begin{tikzpicture}[scale=0.85]
  \node[font=\large] at (0,0) {$\odot$};
  \foreach \r in {0.55,1.05,1.55}
    \draw[omThm,-Stealth] (\r,0) arc (0:350:\r);
  \node[omThm, font=\small] at (45:1.9) {$\vect B$};
  \node[font=\small] at (0,-2.1) {corrente vindo na sua direção};
\end{tikzpicture}\qquad\qquad
\begin{tikzpicture}[scale=0.85]
  \node[font=\large] at (0,0) {$\otimes$};
  \foreach \r in {0.55,1.05,1.55}
    \draw[omThm,-Stealth] (\r,0) arc (360:10:\r);
  \node[omThm, font=\small] at (45:1.9) {$\vect B$};
  \node[font=\small] at (0,-2.1) {corrente se afastando de você};
\end{tikzpicture}

{\small \omterm{def:g11:electric-gravitational-fields:lines}{Linhas de campo} circulares em volta de um fio retilíneo, vistas
de frente ($\odot$: corrente saindo da página; $\otimes$: entrando
nela). Agarre o fio com a mão direita, com o polegar no sentido da
corrente: os dedos se curvam no sentido em que $\vect B$ gira.}
\end{omfigure}

\begin{example}[Um fio é um ímã fraco]\label{ex:g11:magnetic-fields:wire}
A $d = \qty{1.0}{cm}$ de um fio percorrido por $I = \qty{15}{A}$:
$B = \num{2e-7} \times 15 / 0.010 = \qty{0.30}{mT}$ --- seis \omterm{def:g11:electric-gravitational-fields:field}{campos}
terrestres, a partir de uma corrente que já desarmaria um disjuntor
doméstico. \omterm{def:g11:electric-gravitational-fields:field}{Campos} fortes precisam de uma geometria melhor.
\end{example}

\begin{definition}[Bobina e solenoide]\label{def:g11:magnetic-fields:solenoid}
Enrolar o fio concentra o \omterm{def:g11:electric-gravitational-fields:field}{campo} dele. Uma \emph{bobina
chata}\index{bobina} ($N$ espiras num disco) se comporta como um \omterm{def:g11:magnetic-fields:magnet}{ímã}
fino, com uma face norte e a outra sul; um
\emph{solenoide}\index{solenoide} --- um enrolamento cilíndrico longo,
com $n$ espiras por \omterm{def:g10:orders-of-magnitude:unit}{metro} --- se comporta como um \omterm{def:g11:magnetic-fields:magnet}{ímã} em barra.
\end{definition}

\begin{proposition}[Campo dentro de um solenoide]\label{prop:g11:magnetic-fields:solenoidfield}
Dentro de um \omterm{def:g11:magnetic-fields:solenoid}{solenoide} longo, longe das pontas, o \omterm{def:g11:electric-gravitational-fields:field}{campo} é
\emph{uniforme}: paralelo ao eixo e o mesmo em todos os pontos
interiores, de módulo $B = \mu_0\, n\, I$, independente da largura do
tubo. Do lado de fora, o \omterm{def:g11:electric-gravitational-fields:field}{campo} é fraco.
\end{proposition}
\admitted

\begin{remark}
As duas fórmulas são deduzidas com honestidade no volume do primeiro ano
de graduação, a partir da lei que relaciona a circulação de um \omterm{def:g11:electric-gravitational-fields:field}{campo} à
corrente que ele envolve. Repare no que importa: $n$, espiras \emph{por
\omterm{def:g10:orders-of-magnitude:unit}{metro}} --- enrole mais apertado, não mais comprido.
\end{remark}

\begin{method}[A mão direita, duas vezes]\label{met:g11:magnetic-fields:righthand}
\begin{itemize}
  \item \emph{Fio}: polegar = corrente; dedos curvados = sentido de giro
        das linhas.
  \item \emph{\omterm{def:g11:magnetic-fields:solenoid}{Bobina}, \omterm{def:g11:magnetic-fields:solenoid}{solenoide}}: dedos curvados = corrente nas espiras;
        polegar = $\vect B$ interno, apontando para a face norte.
\end{itemize}
\end{method}

\begin{omfigure}
\begin{tikzpicture}[scale=0.9]
  \foreach \x in {0,0.65,...,3.9} {
    \node[font=\small] at (\x,0.85) {$\odot$};
    \node[font=\small] at (\x,-0.85) {$\otimes$};}
  \draw[-Stealth, omThm, thick] (-0.4,0.4) -- (4.3,0.4);
  \draw[-Stealth, omThm, thick] (-0.4,0) -- (4.3,0)
    node[right, font=\small] {$\vect B$};
  \draw[-Stealth, omThm, thick] (-0.4,-0.4) -- (4.3,-0.4);
  \draw[omExo, -Stealth, dashed]
    (4.5,0.7) .. controls (5.9,2.3) and (-2.0,2.3) .. (-0.6,0.7);
\end{tikzpicture}

{\small Um \omterm{def:g11:magnetic-fields:solenoid}{solenoide} cortado ao comprido ($\odot$/$\otimes$: as espiras
atravessando a página): por dentro, um $B = \mu_0 n I$ uniforme ao longo
do eixo; por fora, um \omterm{def:g11:electric-gravitational-fields:field}{campo} de retorno fraco e espalhado (tracejado).}
\end{omfigure}

\begin{example}[Números de um solenoide]\label{ex:g11:magnetic-fields:solenoid}
Com $10$ espiras por centímetro ($n = \qty{1000}{m^{-1}}$) e
$I = \qty{5.0}{A}$:
$B = \num{1.26e-6} \times 1000 \times 5.0 \approx \qty{6.3}{mT}$ --- cem
\omterm{def:g11:electric-gravitational-fields:field}{campos} terrestres, ajustáveis por um botão.
\end{example}

\section{Eletroímãs}

\begin{definition}[Eletroímã]\label{def:g11:magnetic-fields:electromagnet}
Um \emph{eletroímã}\index{eletroímã} é um \omterm{def:g11:magnetic-fields:solenoid}{solenoide} enrolado sobre um
núcleo de ferro doce: o ferro se imanta no \omterm{def:g11:electric-gravitational-fields:field}{campo} da \omterm{def:g11:magnetic-fields:solenoid}{bobina} e o
multiplica, tipicamente por $100$ a $1000$. Ao contrário de um \omterm{def:g11:magnetic-fields:magnet}{ímã}
permanente, ele se desliga junto com a corrente.
\end{definition}

\begin{remark}[Três máquinas]\label{rem:g11:magnetic-fields:machines}
O guindaste do ferro-velho levanta um carro com um \omterm{def:g11:magnetic-fields:electromagnet}{eletroímã} e --- é esse
o ponto --- o \emph{solta} abrindo uma chave. O pequeno \omterm{def:g11:magnetic-fields:electromagnet}{eletroímã} de um
relé puxa um contato até fechá-lo: uma corrente fraca comanda uma forte
(\cref{ch:g11:circuits-and-power}). Um aparelho de ressonância magnética
mantém \qty{1.5}{T} uniformes em volta do paciente com um \omterm{def:g11:magnetic-fields:solenoid}{solenoide}
supercondutor.
\end{remark}

\omterm{def:g10:orders-of-magnitude:oom}{Ordens de grandeza}, da galáxia a uma estrela morta:
\begin{center}
\small
\begin{tabular}{lr}
espaço interestelar & $\sim \qty{e-10}{T}$ \\
\omterm{def:g11:electric-gravitational-fields:field}{campo} na superfície da Terra & \qty{5e-5}{T} \\
\omterm{def:g11:magnetic-fields:magnet}{ímã} de geladeira & $\sim \qty{5e-3}{T}$ \\
\omterm{def:g11:magnetic-fields:magnet}{ímã} de neodímio, no seu polo & $\sim \qty{0.5}{T}$ \\
\omterm{def:g11:magnetic-fields:solenoid}{solenoide} de ressonância magnética & \qty{1.5}{T} \\
superfície de uma estrela de nêutrons & $\sim \qty{e8}{T}$ \\
\end{tabular}
\end{center}

\section{A força de Laplace}

\begin{proposition}[Força de Laplace]\label{prop:g11:magnetic-fields:laplace}
Um fio retilíneo de comprimento $L$, percorrido por uma corrente $I$ num
\omterm{def:g11:electric-gravitational-fields:capacitor}{campo uniforme} $\vect B$ \emph{perpendicular} ao fio, sente a \emph{força
de Laplace}\index{força de Laplace} de módulo $F = I\,L\,B$,
perpendicular tanto ao fio quanto a $\vect B$: mão direita espalmada,
polegar no sentido da corrente, dedos esticados no sentido de $\vect B$
--- a palma empurra no sentido de $\vect F$. Um fio \emph{paralelo} ao
\omterm{def:g11:electric-gravitational-fields:field}{campo} não sente \omterm{def:g10:inertia:force}{força} alguma.
\end{proposition}
\admitted

\begin{remark}[O trilho saltador]\label{rem:g11:magnetic-fields:rail}
O volume do primeiro ano de graduação deduz isso da \omterm{def:g10:inertia:force}{força} magnética sobre
cada carga em movimento. No laboratório: uma barra de cobre repousa sobre
dois trilhos horizontais entre os polos de um \omterm{def:g11:magnetic-fields:magnet}{ímã} (\omterm{def:g11:electric-gravitational-fields:field}{campo} entrando na
página, na figura); feche a chave e a barra dispara ao longo dos trilhos.
Inverta a corrente, ou vire o \omterm{def:g11:magnetic-fields:magnet}{ímã}: ela dispara para o outro lado.
\end{remark}

\begin{omfigure}
\begin{tikzpicture}[scale=0.95]
  \foreach \x in {0.6,1.5,2.9,3.8}
    \foreach \y in {0.35,1.05}
      \node[omThm, font=\small] at (\x,\y) {$\otimes$};
  \node[omThm, font=\small] at (4.35,1.05) {$\vect B$};
  \draw[very thick] (0,0) -- (4.6,0);
  \draw[very thick] (0,1.4) -- (4.6,1.4);
  \draw[omDef, line width=2.5pt] (2.1,-0.12) -- (2.1,1.52);
  \draw[-Stealth, omDef, thick] (0.1,1.6) -- (1.6,1.6)
    node[midway, above, font=\small] {$I$};
  \draw[-Stealth, omDef, thick] (2.32,1.15) -- (2.32,0.45);
  \draw[-Stealth, omDef, thick] (1.6,-0.22) -- (0.1,-0.22);
  \draw[-Stealth, omMeth, very thick] (2.1,0.7) -- (3.4,0.7)
    node[pos=0.85, above, font=\small] {$\vect F$};
\end{tikzpicture}

{\small A experiência dos trilhos: corrente descendo pela barra e \omterm{def:g11:electric-gravitational-fields:field}{campo}
entrando na página ($\otimes$) --- a \omterm{prop:g11:magnetic-fields:laplace}{força de Laplace} $F = ILB$ empurra a
barra ao longo dos trilhos.}
\end{omfigure}

\begin{example}[Números de Laplace]\label{ex:g11:magnetic-fields:laplace}
Uma barra de comprimento $L = \qty{5.0}{cm}$ percorrida por
$I = \qty{8.0}{A}$ num \omterm{def:g11:electric-gravitational-fields:field}{campo} $B = \qty{0.50}{T}$:
$F = 8.0 \times 0.050 \times 0.50 = \qty{0.20}{N}$ --- o \omterm{def:g10:universal-gravitation:weight}{peso} de
\qty{20}{g}, sobre um fio que pesa alguns gramas.
\end{example}

\begin{remark}[O motor de corrente contínua]\label{rem:g11:magnetic-fields:motor}
Dobre o fio numa espira dentro do \omterm{def:g11:electric-gravitational-fields:field}{campo}: os dois lados perpendiculares a
$\vect B$ conduzem correntes opostas, de modo que as \omterm{prop:g11:magnetic-fields:laplace}{forças de Laplace}
sobre eles são opostas --- um binário que faz a espira girar. Meia volta
depois, o binário a giraria de volta, e por isso uma chave rotativa (o
\emph{comutador}\index{comutador}) inverte a corrente a cada meia volta.
Todo ventilador e toda furadeira são essa espira, multiplicada.
\end{remark}

\section{Exercícios}

\begin{exercise}[$\star$]\label{exo:g11:magnetic-fields:1}
Um \omterm{def:g11:magnetic-fields:magnet}{ímã} em barra é serrado ao meio, entre os seus polos. Que polos cada
pedaço carrega? Aproximados de novo, as faces recém-cortadas se atraem ou
se repelem? Algum esquema de corte consegue isolar o polo norte?
\end{exercise}

\begin{exercise}[$\star$]\label{exo:g11:magnetic-fields:2}
Verdadeiro ou falso, com uma razão para cada: (a) duas \omterm{def:g11:magnetic-fields:lines}{linhas de campo
magnético} se cruzam onde o \omterm{def:g11:electric-gravitational-fields:field}{campo} é forte; (b) fora de um \omterm{def:g11:magnetic-fields:magnet}{ímã}, as \omterm{def:g11:electric-gravitational-fields:lines}{linhas
de campo} correm do polo norte para o polo sul; (c) toda \omterm{def:g11:magnetic-fields:lines}{linha de campo
magnético} é fechada; (d) uma agulha se acomoda perpendicular à linha que
passa por ela.
\end{exercise}

\begin{exercise}[$\star$]\label{exo:g11:magnetic-fields:3}
Calcule o \omterm{def:g11:electric-gravitational-fields:field}{campo} a \qty{2.0}{cm} de um fio retilíneo percorrido por
\qty{10}{A}. Quantas vezes os \qty{50}{\micro T} da Terra isso
representa?
\end{exercise}

\begin{exercise}[$\star$]\label{exo:g11:magnetic-fields:4}
Um \omterm{def:g11:magnetic-fields:solenoid}{solenoide} de $800$ espiras enroladas em \qty{40}{cm} conduz
$I = \qty{1.5}{A}$. Calcule $n$ e depois $B$ no interior; em que $B$ se
transforma se a corrente dobrar?
\end{exercise}

\begin{exercise}[$\star$]\label{exo:g11:magnetic-fields:5}
Um trecho de fio de \qty{10}{cm} conduz \qty{5.0}{A} perpendicularmente a
um \omterm{def:g11:electric-gravitational-fields:field}{campo} de \qty{0.20}{T}. Calcule a \omterm{prop:g11:magnetic-fields:laplace}{força de Laplace}; ela é igual ao
\omterm{def:g10:universal-gravitation:weight}{peso} de que massa?
\end{exercise}

\begin{exercise}[$\star\star$]\label{exo:g11:magnetic-fields:6}
No ponto em que uma caminhante está, a \omterm{def:g11:magnetic-fields:earth}{declinação} é de $\ang{8}$ para
oeste. Ela caminha \qty{2.0}{km} ``rumo ao norte pela bússola'': quanto
ela deriva, e para que lado do norte verdadeiro? Que rumo a teria levado
ao norte verdadeiro?
\end{exercise}

\begin{exercise}[$\star\star$]\label{exo:g11:magnetic-fields:7}
A que distância de um fio retilíneo percorrido por \qty{20}{A} o \omterm{def:g11:electric-gravitational-fields:field}{campo}
dele iguala os \qty{50}{\micro T} da Terra? O que isso diz sobre confiar
numa bússola perto de cabos energizados?
\end{exercise}

\begin{exercise}[$\star\star$]\label{exo:g11:magnetic-fields:8}
Projete um \omterm{def:g11:magnetic-fields:solenoid}{solenoide} que produza $B = \qty{10}{mT}$ com
$I = \qty{4.0}{A}$: calcule o $n$ necessário e depois o número de espiras
para uma \omterm{def:g11:magnetic-fields:solenoid}{bobina} de \qty{25}{cm}.
\end{exercise}

\begin{exercise}[$\star\star$]\label{exo:g11:magnetic-fields:9}
Na experiência dos trilhos, a barra (massa \qty{20}{g}, comprimento
$L = \qty{12}{cm}$ entre os trilhos) conduz \qty{10}{A} num \omterm{def:g11:electric-gravitational-fields:field}{campo}
perpendicular de \qty{0.50}{T}. Calcule a \omterm{prop:g11:magnetic-fields:laplace}{força de Laplace} e depois a
aceleração inicial da barra; compare com $g$.
\end{exercise}

\begin{exercise}[$\star\star$]\label{exo:g11:magnetic-fields:10}
Um fio horizontal de massa \qty{5.0}{g} e comprimento \qty{20}{cm} está
num \omterm{def:g11:electric-gravitational-fields:field}{campo} horizontal $B = \qty{0.10}{T}$ perpendicular a ele. Que
corrente faz a \omterm{prop:g11:magnetic-fields:laplace}{força de Laplace} equilibrar o \omterm{def:g10:universal-gravitation:weight}{peso} do fio, de modo que ele
levite? Para onde a \omterm{def:g10:inertia:force}{força} tem de apontar?
\end{exercise}

\begin{exercise}[$\star\star$]\label{exo:g11:magnetic-fields:11}
Dê o sentido da \omterm{prop:g11:magnetic-fields:laplace}{força de Laplace} (ou diga que ela se anula): (a) corrente
para a direita da página, $\vect B$ entrando na página; (b) corrente para
o alto da página, $\vect B$ saindo da página; (c) corrente paralela a
$\vect B$. Use a mão direita espalmada.
\end{exercise}

\begin{exercise}[$\star\star\star$]\label{exo:g11:magnetic-fields:12}
Um fio horizontal corre na direção geográfica norte--sul, \qty{3.0}{cm}
acima de uma bússola; a componente horizontal do \omterm{def:g11:electric-gravitational-fields:field}{campo} terrestre ali vale
\qty{20}{\micro T}. O fio conduz $I = \qty{5.0}{A}$: calcule o \omterm{def:g11:electric-gravitational-fields:field}{campo} dele
na bússola, dê a direção desse \omterm{def:g11:electric-gravitational-fields:field}{campo} e depois o desvio da agulha (os
\omterm{def:g11:electric-gravitational-fields:field}{campos} se somam como vetores).
\end{exercise}

\begin{exercise}[$\star\star\star$]\label{exo:g11:magnetic-fields:13}
O \omterm{def:g11:magnetic-fields:electromagnet}{eletroímã} de um ferro-velho é uma \omterm{def:g11:magnetic-fields:solenoid}{bobina} de $400$ espiras em
\qty{20}{cm}, de \omterm{def:g11:circuits-and-power:resistance}{resistência} \qty{2.0}{\ohm}, ligada a uma fonte de
\qty{12}{V}; o seu núcleo de ferro multiplica por $100$ o \omterm{def:g11:electric-gravitational-fields:field}{campo} da \omterm{def:g11:magnetic-fields:solenoid}{bobina}
nua. Calcule a corrente, o \omterm{def:g11:electric-gravitational-fields:field}{campo} da \omterm{def:g11:magnetic-fields:solenoid}{bobina} nua, o \omterm{def:g11:electric-gravitational-fields:field}{campo} com o núcleo e a
\omterm{def:g10:energy-conservation:power}{potência} que a \omterm{def:g11:magnetic-fields:solenoid}{bobina} dissipa
(\cref{ch:g11:circuits-and-power}). Por que a sucata cai no instante em
que a chave se abre?
\end{exercise}

\begin{exercise}[$\star\star\star$]\label{exo:g11:magnetic-fields:14}
A espira retangular de um motor tem $N = 50$ voltas conduzindo
$I = \qty{2.0}{A}$; os seus dois lados de comprimento \qty{5.0}{cm} são
perpendiculares a um \omterm{def:g11:electric-gravitational-fields:field}{campo} $B = \qty{0.30}{T}$. Calcule a \omterm{def:g10:inertia:force}{força} sobre
cada um desses lados; por que as duas \omterm{def:g10:inertia:force}{forças} são opostas, e o que elas
fazem com a espira? Por que os outros dois lados às vezes não sentem
\omterm{def:g10:inertia:force}{força} alguma? Por que a corrente precisa ser invertida a cada meia volta?
\end{exercise}

\begin{exercise}[$\star\star\star$]\label{exo:g11:magnetic-fields:15}
A que distância de um cabo que conduz \qty{100}{A} seria preciso estar
para que o \omterm{def:g11:electric-gravitational-fields:field}{campo} dele chegasse a \qty{1}{T}? Compare com o raio
milimétrico do cabo e conclua por que \omterm{def:g11:electric-gravitational-fields:field}{campos} da ordem do \omterm{def:g11:magnetic-fields:bfield}{tesla} são
construídos com \omterm{def:g11:magnetic-fields:solenoid}{solenoides}, e não com fios isolados. Depois: o \omterm{def:g11:electric-gravitational-fields:field}{campo} de
\qty{e8}{T} na superfície de uma estrela de nêutrons está quantas
potências de dez acima do da Terra?
\end{exercise}

\section{Problema: a bússola, o guindaste e a ressonância magnética}

\begin{problem}[{Problema de fim de semana --- três ímãs em serviço: uma
bússola cruzando um oceano, um eletroímã levantando um carro e o
solenoide de uma ressonância magnética --- um só campo, cobrindo dez
potências de dez, fazendo três trabalhos}]
\label{pb:g11:magnetic-fields:1}
O mesmo vetor $\vect B$ guia um veleiro com \qty{50}{\micro T}, levanta
sucata de ferro com \qty{1}{T} e forma a imagem de um joelho com
\qty{1.5}{T} --- três máquinas para as três fórmulas deste capítulo.

\medskip
\noindent\textbf{Parte I --- A bússola.}
Na posição do navio, o \omterm{def:g11:electric-gravitational-fields:field}{campo} da Terra tem componente horizontal de
\qty{20}{\micro T} e componente vertical de \qty{44}{\micro T}; a
\omterm{def:g11:magnetic-fields:earth}{declinação} é de $\ang{6}$ para oeste.
\begin{enumerate}
  \item Por que a agulha de uma bússola aponta para o norte (magnético)?
  \item Calcule o módulo do \omterm{def:g11:electric-gravitational-fields:field}{campo} total e o seu ângulo com a horizontal.
  \item Seguindo o ``norte da bússola'' por \qty{30}{km}, a que distância
        a oeste da rota do norte verdadeiro o navio termina?
  \item Uma carga de aço desvia a agulha mais $\ang{4}$ para oeste: mesma
        pergunta.
  \item Perto do polo geomagnético a bússola é inútil: que componente tem
        a culpa, e por quê?
\end{enumerate}

\medskip
\noindent\textbf{Parte II --- O guindaste.}
O \omterm{def:g11:magnetic-fields:electromagnet}{eletroímã} do ferro-velho é um \omterm{def:g11:magnetic-fields:solenoid}{solenoide} de $600$ espiras enroladas em
\qty{30}{cm}, de \omterm{def:g11:circuits-and-power:resistance}{resistência} \qty{1.5}{\ohm}, alimentado por uma fonte de
\qty{12}{V}; o seu núcleo de ferro multiplica por $50$ o \omterm{def:g11:electric-gravitational-fields:field}{campo} da \omterm{def:g11:magnetic-fields:solenoid}{bobina}
nua.
\begin{enumerate}[resume]
  \item Calcule $n$, a densidade de espiras.
  \item Calcule a corrente na \omterm{def:g11:magnetic-fields:solenoid}{bobina}
        (\cref{ch:g11:circuits-and-power}).
  \item Calcule o \omterm{def:g11:electric-gravitational-fields:field}{campo} da \omterm{def:g11:magnetic-fields:solenoid}{bobina} nua, $\mu_0 n I$.
  \item Calcule o \omterm{def:g11:electric-gravitational-fields:field}{campo} com o núcleo; compare com os \qty{0.5}{T} do
        neodímio.
  \item Calcule a \omterm{def:g10:energy-conservation:power}{potência} dissipada na \omterm{def:g11:magnetic-fields:solenoid}{bobina} e depois a \omterm{def:g10:energy-conservation:energy}{energia} de um
        turno de içamento de \qty{10}{\minute}.
  \item Existe \omterm{def:g11:magnetic-fields:magnet}{ímã} permanente com essa intensidade. Por que usar um
        \omterm{def:g11:magnetic-fields:electromagnet}{eletroímã} mesmo assim?
\end{enumerate}

\medskip
\noindent\textbf{Parte III --- O motor do guincho.}
O guincho do guindaste é um motor de corrente contínua: uma espira
retangular de $N = 100$ voltas, com lados de \qty{4.0}{cm}
perpendiculares a $B = \qty{0.25}{T}$, conduzindo $I = \qty{3.0}{A}$.
\begin{enumerate}[resume]
  \item Calcule a \omterm{prop:g11:magnetic-fields:laplace}{força de Laplace} sobre cada feixe de lados
        perpendicular a $\vect B$.
  \item As duas \omterm{def:g10:inertia:force}{forças} são iguais e opostas e, ainda assim, o efeito
        delas não se cancela. O que elas produzem, e por que precisam ser
        \emph{opostas}?
  \item Que \omterm{def:g10:inertia:force}{força} age sobre os outros dois lados, quando eles ficam
        paralelos a $\vect B$?
  \item Depois de meia volta, as mesmas \omterm{def:g10:inertia:force}{forças} desfariam a rotação. O que
        o \omterm{rem:g11:magnetic-fields:motor}{comutador} faz, e quando?
  \item Dê três mudanças de projeto distintas, cada uma capaz de dobrar o
        par de \omterm{def:g10:inertia:force}{forças}.
\end{enumerate}

\medskip
\noindent\textbf{Parte IV --- A ressonância magnética.}
O \omterm{def:g11:magnetic-fields:solenoid}{solenoide} do aparelho mantém \qty{1.5}{T} uniformes com uma corrente
$I = \qty{500}{A}$.
\begin{enumerate}[resume]
  \item Quantas vezes o \omterm{def:g11:electric-gravitational-fields:field}{campo} da Terra são \qty{1.5}{T}?
  \item Calcule a densidade de espiras $n$ necessária \emph{sem} núcleo de
        ferro algum.
  \item Se o enrolamento tivesse \omterm{def:g11:circuits-and-power:resistance}{resistência} \qty{0.20}{\ohm}, que
        \omterm{def:g10:energy-conservation:power}{potência} ele dissiparia? Que propriedade do enrolamento real
        evita isso, ao preço de um frio extremo?
  \item Desfecho: coloque os \omterm{def:g11:electric-gravitational-fields:field}{campos} do navio, do guindaste e da
        ressonância na escada de potências de dez que vai do espaço
        interestelar (\qty{e-10}{T}) a uma estrela de nêutrons
        (\qty{e8}{T}); o que mudou de uma máquina para a outra --- a
        física ou os \omterm{def:g11:circuits-and-power:current}{ampères}?
\end{enumerate}
\end{problem}
